\chapter{Just-in-Time Compilation}\label{ch:just-in-time-compilation}

An ahead-of-time compiler optimizes without knowing inputs, shapes, or hot paths.
A \term{just-in-time compiler} waits until the program runs, watches what actually happens, and then compiles for that reality.
Waiting trades warm-up time for better information.
This chapter shows how to spend that budget wisely.
We build tiered execution from interpreter to baseline to optimizing compiler, with profiling and on-stack replacement connecting the tiers.

\section{Why Compile Late}\label{sec:ch34-why-late}

Static compilers guess at branch directions, call targets, and value types.
Runtime observation replaces those guesses with counts.
A loop that runs a million iterations deserves compilation even if it looked cold statically.
A call site that always targets one method deserves a direct call even if the class hierarchy allows many.
A value that is always a small integer deserves unboxed arithmetic even if the language allows strings.
The just-in-time compiler, usually shortened to JIT, acts on these facts.
It compiles only hot code, specializes it to observed behavior, and keeps a fallback when behavior changes.
Figure~\ref{fig:ch34-tiers} shows the tiered pipeline used through this chapter.
Execution starts in the interpreter with profiling enabled.
Hot methods enter a baseline compiler that produces code quickly with modest optimization.
Hotter methods enter an optimizing compiler that spends more time for better code.
Deoptimization, the subject of the next chapter, returns to a lower tier when assumptions fail.
Table~\ref{tab:ch34-tiers} compares the tiers by compile speed, code quality, and profiling needs.
The table is qualitative because thresholds and budgets are workload-specific.
No tier is best in isolation.
The system wins when most time runs in good code while little time compiles.

\begin{figure}[htbp]
\centering
\begin{tikzpicture}[node distance=1.0cm and 1.4cm]
\node[basic block] (interp) {interpreter\\ profile};
\node[basic block, right=of interp] (base) {baseline JIT\\ fast, modest code};
\node[basic block, right=of base] (opt) {optimizing JIT\\ slow, best code};
\node[basic block, below=of base] (deopt) {deoptimize\\ back to lower tier};
\draw[cfg edge] (interp) -- (base);
\draw[cfg edge] (base) -- (opt);
\draw[cfg edge] (opt) -- (deopt);
\draw[cfg edge] (deopt) -- (interp);
\end{tikzpicture}
\caption{Tiered execution moves hot code upward as counters grow and returns through deoptimization when assumptions fail, so profiling pays only where it matters.}
\label{fig:ch34-tiers}
\end{figure}

\begin{table}[htbp]
\centering
\caption{Tiers trade compile speed against code quality, with profiling connecting each step upward.}
\label{tab:ch34-tiers}
\begin{tabular}{llll}
\toprule
Tier & Compile speed & Code quality & Needs \\
\midrule
Interpreter & None & Slowest & Counters \\
Baseline & Fast & Modest & Hot threshold \\
Optimizing & Slow & Best & Type profiles \\
\bottomrule
\end{tabular}
\end{table}

\section{Profiling and Hotness}\label{sec:ch34-profiling}

A JIT compiles what counters and profiles select.
\term{Method counters} record entries and loop back-edges.
\term{Type profiles} record the observed kinds at selected operations.
\term{Call-site profiles} record the observed targets of each indirect call.
Algorithm~\ref{alg:ch34-tier-up} gives the promotion rule for one method.
It assumes counters increment on entry and on each back-edge and decay or reset on compilation.
It requires separate thresholds for baseline and optimizing tiers.
It returns the tier decision without compiling anything itself.
Counters are cheap but noisy.
A method that runs once with a huge loop looks like steady hot code.
Decay on a timer or on collection separates sustained heat from one burst.
Profiles are richer but costlier.
Recording every type at every operation slows the interpreter visibly.
Production systems profile only chosen bytecodes: arithmetic, property access, and calls.
Each profile site keeps a small bounded cache of recent observations.
When the cache overflows with distinct kinds, the site marks itself polymorphic and stops growing.
That bound keeps profiling memory proportional to program size rather than execution length.

\begin{algorithm}[htbp]
\caption{Tier promotion from counters and profile maturity}\label{alg:ch34-tier-up}
\begin{pseudocode}
\Require method $m$ with entry count $c_e$, back-edge count $c_b$, profile state $p$
\Ensure tier decision for $m$
\State $heat \gets c_e + c_b$
\If{$heat < T_{base}$}
\State \Return stay in interpreter
\EndIf
\If{$heat < T_{opt}$}
\State \Return compile with baseline JIT
\EndIf
\If{$p$ is mature}
\State \Return compile with optimizing JIT
\EndIf
\State \Return compile with baseline JIT and keep profiling
\end{pseudocode}
\end{algorithm}

The returned tier is a desired execution tier, not an instruction to compile again on every counter update.
Compare it with the installed tier and any pending compilation, enqueue at most one request per version, and publish completed code atomically.
Use saturating counters or checked arithmetic so counter overflow cannot turn a hot method cold.
Moving an already active loop into new code additionally requires the on-stack replacement state mapping described below.

\section{Worked Example: Promoting a Hot Loop}\label{sec:ch34-worked}

We apply Algorithm~\ref{alg:ch34-tier-up} to a small summation loop.
Thresholds are $T_{base} = 100$ and $T_{opt} = 1000$.
The method is entered once, so $c_e = 1$.
The loop body runs $1200$ iterations, so $c_b = 1200$.
Heat is $1 + 1200 = 1201$.
After $50$ iterations, heat is $51$, below $T_{base}$, so the method stays interpreted.
After $150$ iterations, heat is $151$, above $T_{base}$ but below $T_{opt}$, so the baseline compiler produces quick code.
After $1200$ iterations, heat is $1201$ and the integer profile at the add is mature, so the optimizing compiler produces specialized code.
The profile matters as much as the count.
Suppose the add observed $1150$ integers and $1$ string before maturity.
The optimizer specializes for integers and guards the string case.
If the string never recurs, the guard never fires and the fast path pays for itself.
If strings arrive steadily, the site is polymorphic and the optimizer should not specialize narrowly.
The listing below shows the specialized shape in schematic intermediate form.
\begin{irlisting}
; specialized loop body, integer profile mature
; guard: tag(addend) == int, else deoptimize
t1 = unchecked_add t0, 1   ; no tag test
t2 = cmp_lt t1, 1200
br t2, loop, exit
\end{irlisting}
Arithmetic checks the decision: $1201 \geq 1000$ promotes, $151 \geq 100$ but $< 1000$ holds at baseline, and $51 < 100$ waits.
Each transition keeps a way back, which the next chapter implements.

\section{On-Stack Replacement}\label{sec:ch34-osr}

Compiling a method is not enough when the program already sits inside a long-running loop.
\term{On-stack replacement} transfers control between tiers mid-activation.
Entry replacement starts a method in compiled code.
Loop replacement jumps from interpreted frames into a compiled loop body, or back out again.
The transfer needs a mapping between interpreter state and compiled state at the chosen point.
Source-level values must match on both sides: locals, stack temporaries, and the program counter.
The compiler emits an \term{OSR map} at each transfer point for this purpose.
\begin{ccnote}[Maps before speed]
An OSR map that drops one local corrupts the program silently.
Test state transfer with exact value checks before measuring any speedup.
\end{ccnote}
Test replacement with a loop that transfers at iteration $500$ of $1000$.
The test records every local before and after the jump and compares checksums over all iterations.
A second test transfers back out at iteration $750$ and finishes in the interpreter.
Both runs must produce the same final sum.
The \enquote{same sum both ways} check catches most mapping defects.

\section{Code Caches and Invalidation}\label{sec:ch34-caches}

Compiled code lives in a managed \term{code cache} with a fixed budget.
Each entry holds machine code, its OSR maps, its profile links, and its dependency list.
When the cache fills, cold entries are evicted and their methods return to interpreted execution.
Eviction must patch or remove every call site that targeted the evicted code.
Indirect calls through a table need only a table update.
Inlined copies need invalidation of each host method that absorbed them.
\term{Invalidation} generalizes eviction to any broken assumption.
Loading a class can add an override that defeats a devirtualized call.
Redefining a method can change semantics the optimizer baked in.
A profile contradiction can prove a specialization wrong.
The system walks dependencies from the changed entity to every affected compilation.
Live activations of discarded code deoptimize at their next safe point, reusing the machinery of the next chapter.
New calls use fresh compilations or the interpreter until heat rebuilds.
Test the cache with a fill-and-evict sequence under fixed profiles.
The test asserts that every evicted method still computes correct results and that hot methods recompile on demand.
A second test loads a new class mid-run and asserts that affected call sites observe the new target.
The \enquote{correct after eviction} rule keeps cache policy from becoming a correctness risk.

\section{Implementation Notes and Complexity}\label{sec:ch34-implementation}

Counter updates cost $O(1)$ per entry and back-edge.
Profile recording costs $O(1)$ per observed site with bounded cache size.
The promotion test itself is $O(1)$.
Compilation cost dominates: baseline compilation is roughly linear in method size, while optimizing compilation includes analyses that are quadratic or worse in pathological cases.
Code caches bound total memory; eviction returns cold methods to interpreted execution.
Production systems compile on background threads, queue by heat, and cap concurrent jobs.
They invalidate compiled code when classes load, methods redefine, or profiles contradict assumptions.
Each invalidation walks a dependency table from the changed entity to affected code.
A common pitfall sets thresholds from folklore instead of measurement.
Too low a threshold compiles cold code and wastes warm-up.
Too high a threshold leaves real heat interpreted.
Another pitfall profiles everything and slows the interpreter severalfold.
Selective profiling at arithmetic, access, and call sites captures most decisions at a fraction of the cost.
We keep thresholds configurable and profiles bounded from the start.
The tiers here feed speculation in the next chapter directly.
Guards compiled now become the deoptimization points handled there.
% TODO(cite): need a primary reference for tiered-compilation threshold policies.

\section{Exercises}\label{sec:ch34-exercises}

\begin{enumerate}
\item Compute the tier decision of Algorithm~\ref{alg:ch34-tier-up} for $c_e = 3$, $c_b = 90$, $T_{base} = 100$ with an immature profile, then repeat with $c_b = 2000$ and a mature profile.\label{ex:ch34-tier-math}
\item Design a call-site profile with three observed targets in counts $98$, $1$, and $1$, state which specialization you would compile, and describe the fallback for the rare targets.\label{ex:ch34-call-profile}
\item Write an on-stack-replacement test for a loop of $1000$ iterations transferring at iteration $400$, listing which values the OSR map must carry and what checksum you would compare.\label{ex:ch34-osr-test}
\end{enumerate}

\section{Further Reading}\label{sec:ch34-reading}

Appel describes compilation pipelines whose ideas carry over to run-time compilation~\cite{appel-1998-modern}.
Morgan relates profiling and code shape to optimization payoffs~\cite{morgan-1998-optimizing}.
Safonov discusses trustworthy compilation when assumptions can change at run time~\cite{safonov-2010-trustworthy}.
The C standard draft defines observable behavior that any replacement must preserve~\cite{wg14-2011-c11}.
